\paragraph{Success criterion.} The criterion of Section~\ref{sec:intro} is not met by gate-model
QAOA: with penalties it learns the constraint rather than the objective (E6, E8b), and with
constraint-preserving mixers it remains well below FeasibleSA and MILP at every depth we simulated
(E8, E8b, E12). Feasible-space PIMC meets a 0.95 bar relative to the best classical method, but it
only ties feasible-space SA, and at equal proposals it is slightly behind (E12b).

\paragraph{What decided the outcome.} Each step that improved a quantum or quantum-inspired method
was a change of \emph{encoding}: certified but gentler penalties (E1, E4, E10), native higher-order
terms instead of quadratisation (E5, E11), and moves or mixers that respect the constraints (E8,
E12). Each of those changes helped the classical methods at least as much. The objective's
submodularity makes greedy strong wherever equity does not bind, and HiGHS certifies every
benchmark instance quickly.

\paragraph{What hardware could still test.} A fair hardware test would submit the constrained
model natively, as a CQM to a hybrid solver, on the E8b ``exactly 3'' city (seed 4) and a family-C
park city (seed 100). A safe-$\Lam$ slack \QUBO{} should never be sent to a QPU, because its gap is
set by $\Lam$ (E6, E11). The repository contains a hook that is inactive without an API token; no
D-Wave result exists.
